% !PRINTED 529
\weberpage{529}
%<<CONT
Discriminanten aller Basen von $\Omega$ dasselbe Vorzeichen haben. Die Formel (10) zeigt auch, dass irgend ein System von $n$ Zahlen $\omega_r$ des Körpers $\Omega$ immer dann eine Basis von $\Omega$ ist, wenn das Determinantenquadrat (7) nicht verschwindet.

Ist $\omega_1, \omega_2, \ldots, \omega_n$ eine Basis von $\Omega$, und bedeuten $c_{r,s}$ rationale Zahlen, so bilden auch die $n$ Zahlen
\[
\omega_r' = c_{r,1}\omega_1 + c_{r,2}\omega_2 + \cdots + c_{r,n}\omega_n \qquad r = 1, 2, \ldots, n
\]
eine Basis von $\Omega$, wenn die Determinante
\[
C = \Sigma \pm c_{1,1}c_{2,2} \ldots c_{n,n}
\]
nicht verschwindet, und es ist
\weq{11}{\varDelta(\omega_1', \omega_2', \ldots, \omega_n') = C^2 \varDelta(\omega_1, \omega_2, \ldots, \omega_n).}

Wenn wir z. B. die Elemente $\omega_1, \omega_2, \ldots, \omega_n$ einer Basis mit rationalen Coëfficienten $c_1, c_2, \ldots, c_n$ multipliciren, deren keiner verschwindet, so erhalten wir eine neue Basis
\[
c_1\omega_1, c_2\omega_2, \ldots, c_n\omega_n.
\]

\begin{center}
\S.\,145.

\gesperrt{Die Minimalbasis und die Körperdiscriminante.}
\end{center}

Nach der zuletzt gemachten Bemerkung verliert eine Basis von $\Omega$ die Eigenschaft, eine Basis zu sein, nicht, wenn man jede ihrer Zahlen mit einer von Null verschiedenen rationalen Zahl multiplicirt. Nun kann man nach \S.\,133,5. jede Zahl durch Multiplication mit einer ganzen rationalen Zahl in eine ganze Zahl verwandeln, und daraus folgt, \gesperrt{dass es Basen von $\Omega$ giebt, deren Elemente lauter ganze Zahlen sind. Die Discriminante einer solchen Basis ist eine ganze rationale Zahl.} Diese ganze rationale Zahl ist von Null verschieden. Sie ändert sich, wenn eine andere ganzzahlige Basis gewählt wird, behält aber für einen bestimmten Körper ein unverändertes Vorzeichen.

Unter all diesen ganzen Zahlen, die als Discriminanten einer ganzzahligen Basis auftreten können, und die alle in quadratischem Verhältniss zu einander stehen, muss nun eine dem absoluten Werthe nach die kleinste sein. Diese kleinste Discriminante bezeichnen wir mit $\varDelta$ und nennen sie die \gesperrt{Grundzahl} oder auch die \gesperrt{Discriminante des Körpers} $\Omega$.
